%% =====================================================================
%% پایان‌نامه کارشناسی — شایان عموزاده — دانشگاه سمنان
%% استاد راهنما: دکتر دری‌گیو (Dr. Dorrigiv)
%% نسخه اصلاح‌شده بر اساس اصلاحیه استاد (۴ بند) — خروجی: کد LaTeX
%% نگاشت اصلاحیه استاد به تغییرات (همه تغییرات با %% FIX علامت‌گذاری شده):
%%   ۱-۱ نام دانشگاه یک‌بار در هر جلد ............ بررسی و تثبیت شد
%%   ۱-۲ Dorrigiv در همه بخش‌های انگلیسی ......... بررسی و تثبیت شد
%%   ۱-۳ بازنویسی چکیده انگلیسی .................. پرداخت نهایی آکادمیک
%%   ۱-۴ حذف صفحه سفید ۳۳ ....................... حذف شد
%%   ۲-۱ شماره‌گذاری فارسی راست‌به‌چپ ............ تثبیت + الگوریتم/کد فارسی
%%   ۲-۲ ارقام فارسی + «جدول» .................... اصلاح سلول‌ها و تیترها
%%   ۲-۳ پانویس اصطلاحات انگلیسی ................. اعمال کامل + حذف انگلیسی میانی
%%   ۳-۱ نمودارها و تصاویر فنی ................... ۸ شکل TikZ + ۲ شماتیک UI
%%   ۳-۲ بسط حجم فصول ........................... شبه‌کد، الگوریتم، تحلیل فنی
%%   ۳-۳ فصل ارزیابی و نتایج ..................... فصل پنجم مستقل (جدید)
%%   ۴-۱ اصلاح ادعای میکروسرویس .................. تثبیت + جدول کنش‌ها
%%   ۴-۲ JWT با HMAC-SHA256 ...................... تثبیت + رابطه راستی‌آزمایی
%% روش کامپایل: XeLaTeX (دو بار اجرا برای فهرست‌ها)، قلم‌های Noto Sans و
%% Noto Sans Arabic. بسته‌های لازم: tikz، algorithm، algorithmic، listings.
%% =====================================================================
\documentclass[11pt, a4paper]{report}
\usepackage[top=2.5cm, bottom=2.5cm, left=2.5cm, right=3.5cm]{geometry}
\usepackage{fontspec}

\usepackage[persian, bidi=basic, provide=*]{babel}

\babelprovide[import, mapdigits, onchar=ids fonts]{persian}
\babelprovide[import, onchar=ids fonts]{english}

% Set default/Latin font to Sans Serif in the main (rm) slot
\babelfont{rm}{Noto Sans}
% Assign a specific font for Persian text
\babelfont[persian]{rm}{Noto Sans Arabic}

\usepackage{enumitem}
\setlist[itemize]{label=\textbullet}
\usepackage{setspace}
\usepackage{titlesec}
\usepackage{tocloft}
\usepackage{amsmath}
\usepackage{booktabs}
\usepackage{array}
\usepackage{tabularx}
\usepackage{float}

%% FIX (3-1): رسم نمودارها با TikZ/PGF طبق دستور استاد (کتابخانه babel برای سازگاری با متن فارسی).
\usepackage{tikz}
\usetikzlibrary{shapes.geometric, arrows.meta, positioning, calc, fit, backgrounds, babel}
%% FIX (3-2): شبه‌کدها و الگوریتم‌ها برای بسط فنی فصول.
\usepackage{algorithm}
\usepackage{algorithmic}
%% FIX (3-2): قطعه‌کدهای کوتاه مستند (شواهد پیاده‌سازی).
\usepackage{listings}
\usepackage{graphicx}

\usepackage[colorlinks=true, linkcolor=black, citecolor=black, urlcolor=blue]{hyperref}

\addto\captionspersian{%
  \renewcommand{\tablename}{جدول}%
  \renewcommand{\listtablename}{فهرست جدول‌ها}%
  \renewcommand{\figurename}{شکل}%
  \renewcommand{\listfigurename}{فهرست شکل‌ها}%
  %% FIX (2-1/3-2): نام فارسی محیط الگوریتم و فهرست آن.
  \renewcommand{\listalgorithmname}{فهرست الگوریتم‌ها}%
}
%% FIX (2-1/3-2): عنوان شناور الگوریتم به فارسی.
\floatname{algorithm}{الگوریتم}

\renewcommand{\thechapter}{\arabic{chapter}}
\renewcommand{\thesection}{\thechapter-\arabic{section}}
\renewcommand{\thesubsection}{\thesection-\arabic{subsection}}
\renewcommand{\thefigure}{\thechapter-\arabic{figure}}
\renewcommand{\thetable}{\thechapter-\arabic{table}}
\renewcommand{\theequation}{\thechapter-\arabic{equation}}
%% FIX (2-1): شماره فارسی فصل-ردیف برای الگوریتم‌ها (راست‌به‌چپ مثل بقیه).
\renewcommand{\thealgorithm}{\thechapter-\arabic{algorithm}}

\newcommand{\enfootnote}[1]{\footnote{\foreignlanguage{english}{#1}}}
%% FIX (2-3): شناسه‌های فنی لاتین (نام کنش‌ها و ستون‌ها) با قلم یکدست و جهت چپ‌به‌راست.
\newcommand{\code}[1]{\foreignlanguage{english}{\texttt{#1}}}

%% FIX (3-2): پیکربندی قطعه‌کدها؛ توضیحات فارسی بیرون از کد و در متن می‌آید.
\lstset{
  basicstyle=\ttfamily\scriptsize,
  keywordstyle=\bfseries,
  stringstyle=\ttfamily,
  showstringspaces=false,
  breaklines=true,
  frame=single,
  numbers=left,
  numberstyle=\tiny,
  captionpos=b
}

%% FIX (3-1): سبک‌های مشترک نمودارهای TikZ (سیاه‌وسفید، مناسب چاپ پایان‌نامه).
\tikzset{
  amlakbox/.style={draw, thick, rounded corners=2pt, align=center, font=\small},
  amlaklane/.style={draw, very thick, rounded corners=4pt},
  amlakarrow/.style={->, >=stealth, thick},
  amlakdb/.style={draw, thick, cylinder, shape border rotate=90, aspect=0.35, align=center, font=\small},
  amlakdecision/.style={draw, thick, diamond, aspect=2.2, align=center, font=\small, text width=2.8cm},
  amlaknote/.style={font=\scriptsize, fill=white, inner sep=1.5pt}
}

\titleformat{\chapter}[display]
  {\normalfont\huge\bfseries\centering}{\chaptertitlename\ \thechapter}{20pt}{\Huge}
\titlespacing*{\chapter}{0pt}{3.5cm}{2cm}

\setstretch{1.5}

\begin{document}

\pagenumbering{alph}

\begin{titlepage}
    \centering
    \vspace*{0.5cm}

    \begin{figure}[h]
      \centering
      \framebox{\parbox{0.35\textwidth}{\centering
        \vspace{1.2cm}
        \textbf{آرم دانشگاه سمنان} \\
        \small\textit{(Semnan\_University\_Logo\_Persian.png)}
        \vspace{1.2cm}
      }}
    \end{figure}
    \vspace{0.3cm}
    %% FIX (1-1): نام «دانشگاه سمنان» دقیقاً یک‌بار در صفحه عنوان فارسی درج شده است.
    {\LARGE \bfseries دانشگاه سمنان \par}
    \vspace{0.5cm}
    {\large \bfseries دانشکده مهندسی برق و کامپیوتر \par}
    \vspace{1.5cm}
    {\normalsize \bfseries پایان‌نامه کارشناسی مهندسی کامپیوتر -- گرایش نرم‌افزار \par}
    \vspace{1.8cm}
    {\huge \bfseries طراحی و پیاده‌سازی سامانه هوشمند مدیریت املاک \par}
    \vspace{0.4cm}
    {\large \bfseries مبتنی بر هوش مصنوعی و معماری مونولیتیک ماژولار \par}
    \vspace{2.2cm}
    {\normalsize \bfseries نگارش: \par}
    {\Large \bfseries شایان عموزاده \par}
    \vspace{1.5cm}
    {\normalsize \bfseries استاد راهنما: \par}
    %% FIX (نگارش): یکدست‌سازی نیم‌فاصله در نام خانوادگی استاد راهنما.
    {\Large \bfseries دکتر دری‌گیو \par}
    \vfill
    {\normalsize \bfseries شهریور ۱۴۰۵ \par}
\end{titlepage}

\clearpage
\thispagestyle{empty}
\mbox{}

\clearpage
\thispagestyle{empty}
\vspace*{7cm}
\begin{center}
    {\Huge \bfseries بسم الله الرحمن الرحیم}
\end{center}

\clearpage
\begin{titlepage}
    \centering
    \vspace*{0.5cm}

    \begin{figure}[h]
      \centering
      \framebox{\parbox{0.35\textwidth}{\centering
        \vspace{1.2cm}
        \textbf{آرم دانشگاه سمنان} \\
        \small\textit{(Semnan\_University\_Logo\_Persian.png)}
        \vspace{1.2cm}
      }}
    \end{figure}
    \vspace{0.3cm}
    %% FIX (1-1): نام «دانشگاه سمنان» دقیقاً یک‌بار در صفحه عنوان داخلی درج شده است.
    {\LARGE \bfseries دانشگاه سمنان \par}
    \vspace{0.5cm}
    {\large \bfseries دانشکده مهندسی برق و کامپیوتر \par}
    \vspace{1.5cm}
    {\normalsize \bfseries پایان‌نامه کارشناسی مهندسی کامپیوتر -- گرایش نرم‌افزار \par}
    \vspace{1.8cm}
    {\huge \bfseries طراحی و پیاده‌سازی سامانه هوشمند مدیریت املاک \par}
    \vspace{0.4cm}
    {\large \bfseries مبتنی بر هوش مصنوعی و معماری مونولیتیک ماژولار \par}
    \vspace{2.2cm}
    {\normalsize \bfseries نگارش: \par}
    {\Large \bfseries شایان عموزاده \par}
    \vspace{1.5cm}
    {\normalsize \bfseries استاد راهنما: \par}
    %% FIX (نگارش): یکدست‌سازی نیم‌فاصله در نام خانوادگی استاد راهنما.
    {\Large \bfseries دکتر دری‌گیو \par}
    \vfill
    {\normalsize \bfseries شهریور ۱۴۰۵ \par}
\end{titlepage}

\clearpage
\thispagestyle{empty}
\vspace*{1.5cm}
\begin{center}
    {\Large \bfseries تأییدیه اصالت رساله/پایان‌نامه \par}
\end{center}
\vspace{1.2cm}
اینجانب \textbf{شایان عموزاده} بدین‌وسیله اظهار می‌دارم که محتوای علمی این نوشتار با عنوان \textbf{«طراحی و پیاده‌سازی سامانه هوشمند مدیریت املاک مبتنی بر هوش مصنوعی و معماری مونولیتیک ماژولار»} که به عنوان پایان‌نامه کارشناسی مهندسی کامپیوتر به دانشکده مهندسی برق و کامپیوتر دانشگاه سمنان ارائه شده، دارای اصالت پژوهشی بوده و حاصل فعالیت علمی اینجانب است.

اینجانب می‌دانم که اگر خلاف ادعای بالا در هر زمان محرز شود، کلیه حقوق مترتب بر این نوشتار از اینجانب سلب شده و مراتب قانونی مرتبط با آن نیز از طرف مراجع ذی‌ربط قابل پیگیری خواهد بود.
\vspace{2.5cm}
\begin{flushleft}
    \textbf{نام و نام خانوادگی:} شایان عموزاده \\
    \textbf{امضا و تاریخ:} .........................
\end{flushleft}

\clearpage
\thispagestyle{empty}
\begin{center}
    {\large به نام خدا \par}
    {\large دانشگاه سمنان \par}
    {\large دانشکده مهندسی برق و کامپیوتر \par}
    \vspace{0.8cm}
    {\Large \bfseries تأیید دفاع از پایان‌نامه کارشناسی \par}
\end{center}
\vspace{0.8cm}
پایان‌نامه آقای \textbf{شایان عموزاده} برای اخذ درجه کارشناسی مهندسی کامپیوتر با عنوان:
\vspace{0.4cm}
\begin{center}
    {\Large \bfseries طراحی و پیاده‌سازی سامانه هوشمند مدیریت املاک مبتنی بر هوش مصنوعی و معماری مونولیتیک ماژولار \par}
\end{center}
\vspace{0.4cm}
در تاریخ ۲۵ شهریور ۱۴۰۵ دفاع شد و مورد تأیید قرار گرفت.
\vspace{1.5cm}

\noindent \textbf{تأیید کنندگان:}
\vspace{0.8cm}

\noindent ۱) استاد محترم داور: ................................ \hfill امضا \par
\vspace{0.8cm}
%% FIX (نگارش): یکدست‌سازی نیم‌فاصله در نام خانوادگی استاد راهنما.
\noindent ۲) استاد محترم راهنما: دکتر دری‌گیو \hfill امضا \par
\vspace{0.8cm}
\noindent ۳) مدیر محترم گروه مهندسی کامپیوتر: ................................ \hfill امضا \par

\clearpage
\thispagestyle{empty}
\vspace*{7cm}
\begin{center}
    {\Large تقدیم به پدر و مادر عزیزم که پشتیبان و راهنمای همیشگی من در مسیر علم و زندگی بودند.}
\end{center}

\clearpage
\thispagestyle{empty}
\vspace*{2.5cm}
\begin{center}
    {\Large \bfseries سپاس‌گزاری}
\end{center}
\vspace{1cm}
%% FIX (نگارش): یکدست‌سازی نیم‌فاصله در نام خانوادگی استاد راهنما.
بر خود واجب می‌دانم از استاد راهنمای گرانقدرم، جناب آقای \textbf{دکتر دری‌گیو}، که با سعه صدر، راهنمایی‌های ارزنده مهندسی و صرف وقت بی‌دریغ، بنده را در انجام این پژوهش کاربردی یاری فرمودند، کمال تشکر و قدردانی را به عمل آورم. همچنین از تمامی اساتید فرهیخته دانشکده مهندسی برق و کامپیوتر دانشگاه سمنان صمیمانه سپاسگزارم.

\clearpage
\pagenumbering{arabic}
\setcounter{page}{1}
\addcontentsline{toc}{chapter}{چکیده}
\begin{center}
    {\Large \bfseries چکیده}
\end{center}
\vspace{0.5cm}
صنعت املاک و مستغلات در عصر تحول دیجیتال نیازمند بهره‌گیری از زیرساخت‌های نرم‌افزاری نوین جهت سرعت‌بخشی به فرآیندهای فایلینگ، افزایش تعامل با مشتریان و تأمین سطح بالایی از محرمانگی داده‌ها است. سیستم‌های سنتی کاغذی و پایگاه‌های داده محلی قدیمی، علاوه بر اتلاف زمان مشاوران، امنیت اطلاعات تجاری آژانس‌ها را با تهدید روبرو می‌سازند. هدف این پایان‌نامه، تحلیل، طراحی و پیاده‌سازی یک سامانه جامع و تحت وب برای مدیریت هوشمند املاک است که بر مبنای معماری مونولیتیک ماژولار در بستر کلاینت-سرور پیاده‌سازی شده و از پاردایم برنامه‌های وب پیش‌رونده\enfootnote{Progressive Web Application (PWA)} بهره می‌جوید.

در این سامانه، ماژول‌های نوین هوش مصنوعی شامل استخراج خودکار مشخصات فایل با مدل زبانی بزرگ، تبدیل گفتار به متن بر پایه وب‌سرویس‌های مدل ویسپر\enfootnote{Whisper ASR Model}، موقعیت‌یابی برخط روی سرویس نقشه نشان و تولید خودکار کاتالوگ اسناد توسعه یافته‌اند. از منظر امنیت معماری نرم‌افزار، با بازنگری اساسی روش‌های منسوخ مبتنی بر نشست و هش‌های ضعیف، مکانیزم احراز هویت بدون حالت مطابق با استاندارد توکن وب جیسون\enfootnote{JSON Web Token (JWT - RFC 7519)} و الگوریتم رمزنگاری اچ‌مک\enfootnote{HMAC-SHA256} پیاده‌سازی شد. همچنین چالش‌های عملیاتی نظیر تناقض نمک توکن در زمان تمدید لایسنس و خطاهای محاسبات زمانی ناشی از اختلاف منطقه زمانی سرور با زمان رسمی کشور شناسایی و با راه‌حل‌های مهندسی مهار گردیدند. ارزیابی‌های عملکردی نشان می‌دهد که زمان استخراج مشخصات از کلام مشاور به کمتر از ۳ ثانیه رسیده و سیستم در برابر حملات تزریق کد و دسترسی غیرمجاز پایداری بالایی دارد.

\vspace{0.8cm}
%% FIX (2-3): حذف واژگان انگلیسی داخل متن فارسی کلیدواژه‌ها؛ معادل‌ها در پانویس چکیده آمده است.
\noindent \textbf{کلیدواژگان:} سامانه مدیریت املاک، هوش مصنوعی، توکن وب جیسون، تبدیل گفتار به متن، معماری مونولیتیک ماژولار، وب‌اپلیکیشن پیش‌رونده.

\clearpage
\tableofcontents
\clearpage
\listoffigures
\clearpage
\listoftables
\clearpage
%% FIX (3-2): فهرست الگوریتم‌های اضافه‌شده به پایان‌نامه.
\listofalgorithms
\clearpage

\chapter{مقدمه}

\section{پیشگفتار و بیان مسئله}
مدیریت دقیق اطلاعات، تطبیق سریع تقاضای مشتریان با فایل‌های بایگانی‌شده و جلوگیری از اشتراک‌گذاری ناخواسته داده‌های محرمانه مالکین، از چالش‌های بنیادین مدیران و مشاوران املاک به شمار می‌رود. در دهه‌های پیشین، استفاده از روش‌های فیزیکی مانند دفاتر سنتی املاک رواج داشت که این امر علاوه بر اتلاف بخش عمده‌ای از زمان کاری مشاوران، موجب از دست رفتن فرصت‌های کلیدی در بازار پویا و رقابتی مسکن می‌شد.

با نفوذ رایانه‌ها به کسب‌وکارها، ابزارهای مدیریت املاک متولد شدند تا فرآیندهای فایلینگ را الکترونیکی کنند. با این وجود، ورود به عصر ارتباطات دیجیتال نیاز به بستر مقیاس‌پذیر، قابل دسترس در تمامی دستگاه‌ها و مجهز به مکانیزم‌های اتوماسیون هوشمند را به عنوان ضرورتی انکارناپذیر مطرح ساخته است.

\section{سیر تکامل سامانه‌های مدیریت املاک}
فناوری‌های فایلینگ املاک در دهه‌های اخیر سه نسل اساسی را پشت سر گذاشته‌اند:
\begin{enumerate}
    \item \textbf{نسل اول (دفاتر کاغذی و بایگانی دستی):} کاملاً سنتی، بدون کمترین قابلیت جستجوی ساختاریافته و دارای ریسک بسیار بالای مفقودی یا سرقت اطلاعات مشتریان.
    \item \textbf{نسل دوم (نرم‌افزارهای دسکتاپ و پایگاه‌های داده محلی):} نرم‌افزارهای ویندوزی مبتنی بر پایگاه‌های داده محلی که امکان جستجوی اولیه را فراهم می‌ساختند، اما دسترسی به اطلاعات صرفاً به محیط فیزیکی دفتر کار محدود بود و امکان هماهنگی خارج از دفتر وجود نداشت.
    \item \textbf{نسل سوم (سامانه‌های ابری پایه):} وب‌سایت‌هایی که دسترسی از راه دور را محقق کردند، اما همچنان وابسته به ورود دستی ده‌ها فیلد اطلاعاتی از طریق صفحه‌کلید بوده و از قابلیت‌های پردازش صوت، تطبیق هوشمند و امنیت پیشرفته توکن بی‌بهره ماندند.
\end{enumerate}

\section{چالش‌ها و کاستی‌های سامانه‌های موجود}
بررسی میدانی راهکارهای فعال در بازار مسکن نشان‌دهنده کاستی‌های فنی مهمی است:
\begin{itemize}
    \item \textbf{ورود داده‌های زمان‌بر:} الزام مشاور به تایپ طولانی مشخصات ملک در هنگام بازدیدهای میدانی موجب کندی در ثبت و کاهش دقت داده‌ها می‌شود.
    \item \textbf{ضعف امنیتی در کنترل نشست‌ها:} اتکای سیستم‌های قدیمی به نشست‌های منسوخ و عدم امضای دیجیتال داده‌ها، سیستم را در برابر سرقت هویت و دستکاری درخواست‌ها آسیب‌پذیر می‌سازد.
    \item \textbf{فقدان انعطاف‌پذیری کاربری:} عدم بهینه‌سازی سامانه‌ها در قالب برنامه‌های کاربردی تک‌صفحه‌ای\enfootnote{Single Page Application (SPA)}، تجربه کاربری کندی را در گوشی‌های هوشمند رقم می‌زند.
\end{itemize}

\section{اهداف و ضرورت انجام پروژه}
این پژوهش با هدف طراحی و پیاده‌سازی سامانه‌ای جامع، امن و هوشمند انجام پذیرفته است تا زمان ثبت داده‌های ملکی را از چند دقیقه به چند ثانیه تعامل صوتی تبدیل کند. این سامانه با ارتقای هسته احراز هویت به استاندارد نوین توکن و ادغام با وب‌سرویس‌های بومی نقشه، پاسخی مهندسی به نیازهای صنف مشاوران املاک ارائه می‌دهد.

%% FIX (3-2): اهداف به‌صورت سنجه‌پذیر بسط داده شد تا مبنای فصل ارزیابی باشند.
اهداف این پژوهش در قالب سنجه‌های قابل ارزیابی زیر صورت‌بندی می‌شود: (الف) کاهش زمان ثبت یک فایل ملکی از مسیر صوتی به کمتر از چند ثانیه؛ (ب) دستیابی به دقت بالای استخراج فیلدهای کلیدی (قیمت، متراژ، موقعیت و امکانات) از کلام فارسی مشاور؛ (پ) استقرار احراز هویت بدون حالت و مقاوم در برابر دستکاری توکن؛ (ت) ارائه تجربه کاربری یکپارچه در دسکتاپ و موبایل در قالب وب‌اپلیکیشن پیش‌رونده. تحقق هر یک از این اهداف در فصل پنجم با سناریوهای آزمون مشخص ارزیابی می‌شود.

\section{ساختار کلی پایان‌نامه}
%% FIX (3-3): به‌روزرسانی نقشه راه با افزودن فصل مستقل ارزیابی (فصل پنجم) و انتقال جمع‌بندی به فصل ششم.
این پایان‌نامه در شش فصل تدوین گردیده است:
\begin{itemize}
    \item \textbf{فصل اول:} تعاریف پایه، تاریخچه، بیان مسئله و ضرورت طرح را بررسی می‌کند.
    \item \textbf{فصل دوم:} مبانی تئوری معماری نرم‌افزار، استانداردهای احراز هویت و مفاهیم هوش مصنوعی را مرور می‌نماید.
    \item \textbf{فصل سوم:} به تشریح معماری سامانه، ساختار پایگاه داده، نمودارها و مدل‌های دسترسی می‌پردازد.
    \item \textbf{فصل چهارم:} جزئیات پیاده‌سازی کلاینت، توابع سرور و چالش‌های پیچیده برنامه‌نویسی را تبیین می‌کند.
    \item \textbf{فصل پنجم:} پروتکل ارزیابی، سناریوهای آزمون، نتایج کارایی هوش مصنوعی و آزمون‌های امنیتی را تحلیل می‌نماید.
    \item \textbf{فصل ششم:} نتایج به‌دست‌آمده را جمع‌بندی نموده و پیشنهادهایی برای توسعه‌های آتی ارائه می‌دهد.
\end{itemize}

\chapter{مبانی تئوری و نیازمندی‌های سیستم}

\section{مقدمه}
در این فصل مبانی مهندسی، معماری‌های انتخاب‌شده، استانداردهای امنیتی ارتباطی و الگوریتم‌های هوش مصنوعی که چارچوب نظری سامانه را تشکیل می‌دهند به تفکیک تبیین می‌گردند.

\section{نیازمندی‌های مهندسی نرم‌افزار}
تحلیل دقیق نیازمندی‌ها نقطه آغازین هر سامانه سازمانی است. برای این سیستم، نیازمندی‌ها در دو بخش مدون شدند:

\subsection{نیازمندی‌های عملکردی}
\begin{itemize}
    \item قابلیت احراز هویت امن با جداسازی کامل سه نقش کاربری مدیر، مشاور و مهمان بر اساس کنترل دسترسی مبتنی بر نقش\enfootnote{Role-Based Access Control (RBAC)}.
    \item امکان ضبط مستقیم صدا از طریق مرورگر، انتقال امن به سرور و استخراج فیلدهای قیمتی و کاربری.
    \item ثبت و جانمایی نقطه دقیق ملک بر روی نقشه تعاملی و محاسبه فاصله.
    \item تولید خروجی کاتالوگ شکیل با فرمت اسناد قابل‌حمل\enfootnote{Portable Document Format (PDF)} برای ارائه به متقاضیان بدون افشای اطلاعات مالک.
\end{itemize}

\subsection{نیازمندی‌های غیرعملکردی}
\begin{itemize}
    \item \textbf{پاسخ‌دهی سریع:} پردازش کوئری‌ها در کمتر از ۱ ثانیه در شرایط بار ترافیکی عادی.
    \item \textbf{امنیت و محرمانگی:} ایمن‌سازی نقطه پایانی\enfootnote{Endpoint} سامانه در برابر حملات تزریق اسکریپت و کدهای مخرب.
    \item \textbf{پایداری و در دسترس بودن:} پایداری عملیاتی سامانه با اتکا به فایلینگ سبک و مکانیزم‌های کش مرورگر.
\end{itemize}

\section{تحلیل معماری: مونولیتیک ماژولار در برابر میکروسرویس}
در ابتدای فاز تحلیل، تفکیک سامانه به میکروسرویس‌های توزیع‌شده با کانتینرهای مستقل بررسی شد. اگرچه معماری میکروسرویس مقیاس‌پذیری بالایی در ابعاد سازمانی کلان فراهم می‌کند، اما سربار ارتباطات بین‌سرویسی از طریق شبکه، پیچیدگی استقرار، هزینه‌های زیرساختی و چالش‌های حفظ تراکنش‌های یکپارچه دیتابیس در پروژه‌های املاک با ابعاد متوسط، گزینه‌ای غیربهینه به شمار می‌رود.

%% FIX (2-3/4-1): انتقال معادل انگلیسی به پانویس و تصریح معماری واقعی سامانه.
بنابراین، \textbf{معماری مونولیتیک ماژولار}\enfootnote{Modular Monolith} در بستر کلاینت-سرور برگزیده شد. در این مدل، لایه رابط کاربری فرانت‌اند به عنوان یک برنامه تک‌صفحه‌ای کاملاً ایزوله توسعه یافته و از طریق واسط‌های برنامه‌نویسی وبی مبتنی بر کنش\enfootnote{Action-Based Web APIs} بر بستر پروتکل انتقال ابرمتن\enfootnote{Hypertext Transfer Protocol (HTTP)} با سرور ارتباط برقرار می‌نماید. هسته سرور بک‌اند به عنوان یک واحد متمرکز اما ماژولاربندی‌شده پیاده‌سازی گشته که درخواست‌ها را از یک نقطه پایانی مرکزی هدایت نموده و منطق پردازشی هر ماژول را در بخش‌های مستقل مدیریت می‌کند.

%% FIX (3-2/4-1): تبیین دقیق تفاوت واسط مبتنی بر کنش با معماری بازنمودی محض.
\subsection{تفاوت واسط مبتنی بر کنش با معماری بازنمودی محض}
الگوی ارتباطی این سامانه از اصول معماری بازنمودی\enfootnote{Representational State Transfer (REST)} مانند عدم حالت‌مندی، قالب‌بندی یکداخت جیسون\enfootnote{JavaScript Object Notation (JSON)} برای تبادل داده و استفاده از کدهای وضعیت استاندارد الهام گرفته است، اما یک پیاده‌سازی بازنمودی محض با نشانی‌دهی منبع‌محور نیست؛ بلکه عملیات از یک نشانی واحد و با پارامتر کنش\enfootnote{Action} تفکیک می‌شوند. این انتخاب آگاهانه با محدودیت‌های میزبانی اشتراکی، سادگی استقرار اتمی و یکپارچگی تراکنش‌های پایگاه داده در یک سامانه املاک با ابعاد متوسط توجیه می‌شود و جزئیات آن در \cite{rest} ریشه نظری دارد. جدول \ref{tab:api_actions} فهرست کنش‌های پیاده‌سازی‌شده در نقطه پایانی مرکزی، روش فراخوانی و سطح دسترسی هر یک را مستند می‌کند.

\begin{table}[htbp]
\centering
\caption{فهرست کنش‌های نقطه پایانی مرکزی و سطح دسترسی آن‌ها}
\label{tab:api_actions}
\renewcommand{\arraystretch}{1.5}
\begin{tabular}{|c|c|c|p{5.5cm}|}
\hline
\textbf{نام کنش} & \textbf{روش} & \textbf{دسترسی مجاز} & \textbf{کاربرد} \\ \hline
\code{loginManager} & \code{POST} & عمومی & ورود مدیر آژانس و صدور توکن مدیریتی. \\ \hline
\code{loginConsultant} & \code{POST} & عمومی & ورود مشاور آژانس و صدور توکن مشاور. \\ \hline
\code{getData} & \code{GET} & عمومی و مهمان & دریافت یکجای فایل‌ها، تقاضاها و اعضا با قلمرو آژانس. \\ \hline
\code{saveProperty} & \code{POST} & مدیر و مشاور & ثبت یا به‌روزرسانی فایل ملکی. \\ \hline
\code{deleteProperty} & \code{POST} & فقط مدیر & حذف فایل ملکی از آژانس. \\ \hline
\code{saveDemand} & \code{POST} & مدیر و مشاور & ثبت یا به‌روزرسانی تقاضای مشتری. \\ \hline
\code{deleteDemand} & \code{POST} & فقط مدیر & حذف تقاضا. \\ \hline
\code{saveMember} & \code{POST} & فقط مدیر & افزودن یا ویرایش عضو و مشاور. \\ \hline
\code{getNotes} & \code{POST} & کاربر واردشده & خواندن یادداشت‌های شخصی کاربر. \\ \hline
\code{saveNotes} & \code{POST} & کاربر واردشده & ذخیره یادداشت‌های شخصی کاربر. \\ \hline
\code{transcribe\_audio} & \code{POST} & مدیر و مشاور & تبدیل فایل صوتی به متن خام. \\ \hline
\code{jarvisProcess} & \code{POST} & مدیر و مشاور & استخراج هوشمند فیلدها از متن. \\ \hline
\end{tabular}
\end{table}

پاسخ موفق هر کنش در قالب یک پاکت یکداخت با کلید پاسخ و پاسخ خطا با کلید خطا و پیام فارسی کاربرپسند بازگردانده می‌شود؛ نقض احراز هویت کد ۴۰۳ و عبور از سقف نرخ درخواست کد ۴۲۹ تولید می‌کند تا کلاینت بتواند رفتار متناسب (مانند خروج خودکار) را اتخاذ کند.

%% FIX (2-3): حذف سرواژه لاتین از تیتر بخش.
\section{مبانی احراز هویت مبتنی بر توکن}
%% FIX (2-3): معادل انگلیسی به پانویس منتقل شد تا متنی فارسی از واژه لاتین تهی شود.
در روش‌های سنتی مدیریت نشست در سرور، داده‌های ورود در حافظه سرور یا دیسک ذخیره می‌شد که با خاموش شدن سرور یا توزیع بار میان چند سرور، وضعیت نشست‌ها مختل می‌گردید. روش مدرن احراز هویت مبتنی بر توکن وب جیسون\enfootnote{JSON Web Token (JWT)} با هدف برقراری ارتباط بدون حالت\enfootnote{Stateless} تدوین شده است.

%% FIX (2-3): پانویس اجزای سه‌گانه توکن در اولین کاربرد.
توکن‌ها حاوی سه قسمت سرآیند\enfootnote{Header}، بارداده\enfootnote{Payload} و امضای رمزنگاری‌شده\enfootnote{Signature} هستند که به وسیله کاراکتر نقطه از هم جدا می‌شوند. سرور با استفاده از یک کلید محرمانه اختصاصی، امضای توکن را ایجاد می‌کند و در هر تراکنش، صحت امضا را بدون نیاز به کوئری سنگین بررسی می‌نماید. مشخصات کامل این استاندارد در \cite{jwt} و تحلیل الگوهای احراز هویت بدون حالت در کاربردهای وب در \cite{auth} آمده است.

\section{تکنولوژی‌های صوتی و مدل‌های زبانی هوش مصنوعی}
استفاده از هوش مصنوعی در این پژوهش مبتنی بر دو گام تکمیلی است:
\begin{itemize}
    %% FIX (2-3): سرواژه‌های انگلیسی به پانویس منتقل شدند.
    \item \textbf{سیستم‌های بازشناسی خودکار گفتار}\enfootnote{Automatic Speech Recognition (ASR)}: مدل‌های پیشرفته شبکه‌های عصبی عمیق مانند ویسپر بر پایه معماری ترانسفورمر\enfootnote{Transformer} توسعه یافته‌اند و توانایی بالایی در فیلتر کردن نویز محیط و درک لهجه‌های فارسی دارند.
    \item \textbf{مدل‌های زبانی بزرگ}\enfootnote{Large Language Models (LLM)}: مدل‌های پردازش زبان طبیعی قادرند با مهندسی دستورالعمل ورودی، متن خام بدون ساختار را فهمیده و داده‌های حیاتی را به عنوان پارامترهای فرم استخراج نمایند.
\end{itemize}

%% FIX (3-2): بسط مبانی نظری صوت و زبان با ارجاع به منابع معتبر.
مدل ویسپر با آموزش ضعیف‌نظارت‌شده در مقیاس بسیار بزرگ روی صدها هزار ساعت گفتار چندزبانه توسعه یافته و جزئیات آن در \cite{whisper} گزارش شده است؛ معماری ترانسفورمر که ستون فقرات هر دو خانواده مدل فوق است نیز در \cite{transformer} معرفی گردیده. در این سامانه، خروجی متن خام مرحله بازشناسی گفتار به‌عنوان ورودی مرحله استخراج موجودیت‌ها عمل می‌کند و هر دو مرحله در فصل چهارم با جزئیات مهندسی تشریح می‌شوند.

\chapter{طراحی معماری و ساختار پایگاه داده}

\section{مقدمه}
در این فصل طراحی فنی سیستم، ساختار پایگاه داده، دیاگرام‌های ارتباطی موجودیت‌ها، مدل کنترل دسترسی کاربران و شیوه رمزنگاری توکن‌های امنیتی تشریح می‌گردد.

\section{بلوک دیاگرام کلی معماری سامانه}
ارتباط میان ماژول‌های کلاینت، سرور متمرکز و وب‌سرویس‌های جانبی در شکل \ref{fig:architecture} ترسیم شده است.

%% FIX (3-1): جایگزینی قاب خالی با بلوک‌دیاگرام واقعی TikZ (کلاینت، سرور، پایگاه داده و سرویس‌های ابری).
\begin{figure}[htbp]
  \centering
  \begin{tikzpicture}[x=1cm, y=1cm]
    % --- lanes ---
    \node[amlaklane, minimum width=4.6cm, minimum height=7.2cm] (laneC) at (12.0,-4.2) {};
    \node[amlaklane, minimum width=4.4cm, minimum height=7.2cm] (laneS) at (7.1,-4.2) {};
    \node[amlaklane, minimum width=4.4cm, minimum height=7.2cm] (laneE) at (2.3,-4.2) {};
    \node[font=\bfseries\small] at (12.0,-0.95) {کاربر و مرورگر (کلاینت)};
    \node[font=\bfseries\small] at (7.1,-0.95) {سرور متمرکز مونولیتیک};
    \node[font=\bfseries\small] at (2.3,-0.95) {سرویس‌های ابری بیرونی};
    % --- client boxes ---
    \node[amlakbox, minimum width=3.9cm, minimum height=0.85cm] (c1) at (12.0,-1.95) {داشبورد و فهرست فایل‌ها};
    \node[amlakbox, minimum width=3.9cm, minimum height=0.85cm] (c2) at (12.0,-3.0) {ضبط صوت و دستیار جارویس};
    \node[amlakbox, minimum width=3.9cm, minimum height=0.85cm] (c3) at (12.0,-4.05) {نقشه تعاملی و موقعیت‌یابی};
    \node[amlakbox, minimum width=3.9cm, minimum height=0.85cm] (c4) at (12.0,-5.1) {کاتالوگ اسناد در مرورگر};
    \node[amlakbox, minimum width=3.9cm, minimum height=0.85cm] (c5) at (12.0,-6.15) {مانیفست و سرویس‌وورکر};
    \node[font=\scriptsize] at (12.0,-7.15) {تک‌صفحه‌ای و پیش‌رونده};
    % --- server boxes ---
    \node[amlakbox, minimum width=3.7cm, minimum height=0.85cm] (s1) at (7.1,-1.95) {مسیریاب کنش‌ها};
    \node[amlakbox, minimum width=3.7cm, minimum height=0.85cm] (s2) at (7.1,-3.0) {احراز هویت توکن};
    \node[amlakbox, minimum width=3.7cm, minimum height=0.85cm] (s3) at (7.1,-4.05) {فایل، تقاضا، عضو، یادداشت};
    \node[amlakbox, minimum width=3.7cm, minimum height=0.85cm] (s4) at (7.1,-5.1) {درگاه تبدیل صوت به متن};
    \node[font=\scriptsize] at (7.1,-6.15) {پی‌اچ‌پی متمرکز و ماژولار};
    \node[font=\scriptsize] at (7.1,-6.7) {نقطه پایانی واحد کنش‌محور};
    % --- external boxes ---
    \node[amlakbox, minimum width=3.7cm, minimum height=0.85cm] (e1) at (2.3,-1.95) {خدمت نقشه نشان};
    \node[amlakbox, minimum width=3.7cm, minimum height=0.85cm] (e2) at (2.3,-3.0) {بازشناسی گفتار ویسپر};
    \node[amlakbox, minimum width=3.7cm, minimum height=0.85cm] (e3) at (2.3,-4.05) {مدل زبانی بیرونی};
    \node[font=\scriptsize] at (2.3,-5.1) {فراخوانی امن سمت سرور};
    % --- database ---
    \node[amlakdb, minimum width=3.2cm, minimum height=1.1cm] (db) at (7.1,-8.9) {پایگاه داده رابطه‌ای};
    % --- arrows ---
    \draw[<->, >=stealth, thick] (9.7,-3.0) -- (9.35,-3.0) node[amlaknote, midway] {جیسون امن};
    \draw[<->, >=stealth, thick] (4.9,-3.0) -- (4.55,-3.0) node[amlaknote, midway] {وب‌سرویس};
    \draw[amlakarrow] (7.1,-7.8) -- (7.1,-8.3) node[amlaknote, midway, right] {دسترسی امن};
    \draw[amlakarrow] (10.0,-6.6) .. controls (7.0,-7.6) and (5.0,-6.8) .. (4.5,-5.6)
      node[amlaknote, pos=0.5, below] {کاشی نقشه};
  \end{tikzpicture}
  \caption{بلوک دیاگرام معماری سامانه و جریان تبادل داده‌ها در بستر کلاینت-سرور}
  \label{fig:architecture}
\end{figure}

نکته کلیدی این نگاشت آن است که تولید کاتالوگ اسناد در سمت مرورگر انجام می‌شود و بار پردازشی آن بر سرور تحمیل نمی‌گردد؛ همچنین کلیدهای محرمانه وب‌سرویس‌های هوش مصنوعی صرفاً در سمت سرور نگهداری می‌شوند و هرگز در اختیار کلاینت قرار نمی‌گیرند. تبادل کلاینت و سرور بر بستر امن انتقال ابرمتن\enfootnote{Hypertext Transfer Protocol Secure (HTTPS)} و در قالب جیسون صورت می‌پذیرد.

%% FIX (2-3): حذف سرواژه لاتین از تیتر بخش.
\section{طراحی پایگاه داده و دیاگرام ارتباط موجودیت‌ها}
%% FIX (2-3): پانویس اصطلاح تزریق در اولین کاربرد؛ شناسه فنی در متن فارسی حذف شد.
ساختار پایگاه داده سیستم برای پوشش مشخصات متغیر املاک به صورت ساختاریافته طراحی شده و از طریق رابط اشیاء داده پی‌اچ‌پی\enfootnote{PHP Data Objects (PDO)} به همراه گزاره‌های آماده هدایت می‌شود تا ریسک تزریق به پایگاه داده\enfootnote{SQL Injection} به صفر برسد. روابط میان موجودیت‌های سیستم در شکل \ref{fig:erd} مشخص است.

%% FIX (3-1): جایگزینی قاب خالی با نمودار ارتباط موجودیت‌های واقعی با TikZ.
\begin{figure}[htbp]
  \centering
  \begin{tikzpicture}[x=1cm, y=1cm]
    \tikzset{ent/.style={draw, thick, rounded corners=2pt, align=center, text width=4.1cm, font=\small, inner sep=4pt}}
    \node[ent] (ag) at (7.25,-1.5) {\textbf{آژانس‌ها} \\ {\footnotesize شناسه، نام، شهر، تلفن، نام مدیر، انقضای لایسنس، نوع پلن}};
    \node[ent] (mb) at (2.6,-5.0) {\textbf{اعضا و مشاوران} \\ {\footnotesize شناسه، شناسه آژانس، نام، نقش، وضعیت، گذرواژه، آخرین بازدید}};
    \node[ent] (pr) at (11.9,-5.0) {\textbf{فایل‌های ملکی} \\ {\footnotesize شناسه، آژانس، نویسنده، شهر، موقعیت، مختصات، متراژ، امکانات، نوع واگذاری، تصاویر}};
    \node[ent] (nt) at (2.6,-8.5) {\textbf{یادداشت‌های شخصی} \\ {\footnotesize شناسه، شناسه آژانس، شناسه عضو، متن، زمان به‌روزرسانی}};
    \node[ent] (dm) at (11.9,-8.5) {\textbf{تقاضاهای مشتریان} \\ {\footnotesize شناسه، آژانس، نام مشتری، بودجه، متراژ، پیگیری}};
    \draw[thick] (ag.west) -- (mb.north) node[amlaknote, pos=0.25] {۱} node[amlaknote, pos=0.75] {N};
    \draw[thick] (ag.east) -- (pr.north) node[amlaknote, pos=0.25] {۱} node[amlaknote, pos=0.75] {N};
    \draw[thick] (mb.south) -- (nt.north) node[amlaknote, pos=0.3] {۱} node[amlaknote, pos=0.7] {N};
    \draw[thick] (ag.south east) .. controls (10.5,-4.5) and (11.9,-5.5) .. (dm.north)
      node[amlaknote, pos=0.2] {۱} node[amlaknote, pos=0.8] {N};
  \end{tikzpicture}
  \caption{طرح‌واره دیاگرام ارتباط موجودیت‌های پایگاه داده سامانه}
  \label{fig:erd}
\end{figure}

جدول \ref{tab:db_schema} مشخصات جداول اصلی سیستم را مستند می‌کند.

%% FIX (2-2/4-1): اصلاح نام ستون‌ها مطابق پیاده‌سازی واقعی (نوع کاربری، نویسنده، وضعیت، پرچم‌های نمایش مهمان).
\begin{table}[htbp]
\centering
\caption{مستندات ساختار اصلی جداول پایگاه داده}
\label{tab:db_schema}
\renewcommand{\arraystretch}{1.5}
\begin{tabular}{|p{2.6cm}|p{4.6cm}|p{6.3cm}|}
\hline
\textbf{نام جدول} & \textbf{فیلدهای کلیدی} & \textbf{شرح و کاربرد} \\ \hline
\code{agencies} & \code{id, name, expireAt, plan\_type, adminPin} & ذخیره اطلاعات آژانس، تاریخ انقضای لایسنس، نوع پلن و گذرواژه درهم‌سازی‌شده مدیریت. \\ \hline
\code{members} & \code{id, agencyId, role, name, pin, lastSeen} & اعضای آژانس و مشاوران، سطح دسترسی و زمان برچسب آخرین بازدید جهت پایش حضور آنلاین. \\ \hline
\code{properties} & \code{id, agencyId, authorName, status, city, location, lat, lng, usage\_type, area, dealType, images} & فایل‌های ملکی، نویسنده، وضعیت انتشار، مختصات جغرافیایی، نوع واگذاری و فهرست جیسونی مسیر تصاویر. \\ \hline
\code{demands} & \code{id, agencyId, clientName, budget, area} & تقاضاهای مشتریان متقاضی ملک به همراه بازه قیمتی و تاریخ‌های پیگیری خودکار. \\ \hline
\end{tabular}
\end{table}

ستون تصاویر از نوع متن جیسونی است و صرفاً فهرست مسیر فایل‌های بهینه‌سازی‌شده را نگه می‌دارد؛ این انتخاب به معنای غیررابطه‌ای بودن پایگاه داده نیست و تمامیت ارجاعی از طریق کلیدهای خارجی منطقی در سطح برنامه حفظ می‌شود.

\section{کنترل دسترسی مبتنی بر نقش}
%% FIX (2-3): حذف سرواژه لاتین از تیتر بخش؛ معادل در پانویس فصل دوم آمده است.
در این سامانه سه سطح کاربری پیاده‌سازی شده و مجوزهای مربوط به هر یک به دقت در جدول \ref{tab:rbac} تفکیک شده است.

%% FIX (2-2/4-1): اصلاح غلط املایی «سامane»، حذف لاتین از عنوان جدول و تطبیق سطرها با رفتار واقعی نقش‌ها.
\begin{table}[htbp]
\centering
\caption{سطوح دسترسی و امکانات کاربران در سامانه}
\label{tab:rbac}
\renewcommand{\arraystretch}{1.5}
\begin{tabular}{|p{2.6cm}|p{5.5cm}|p{5.7cm}|}
\hline
\textbf{نقش کاربری} & \textbf{دسترسی‌های مجاز} & \textbf{محدودیت‌ها} \\ \hline
\textbf{مدیر آژانس} & مدیریت تمامی فایل‌ها، تایید یا مسدودسازی مشاورین، مشاهده داشبورد آماری و تمدید لایسنس. & دسترسی کامل به داده‌های درون آژانس خود. \\ \hline
\textbf{مشاور} & مشاهده و ویرایش فایل‌های آژانس خود، ثبت فایل ملکی، دسترسی به هوش مصنوعی و ثبت تقاضا. & عدم حذف فایل و تقاضا (ویژه مدیر)، عدم دسترسی به آمار مدیر و یادداشت‌های خصوصی سایر همکاران. \\ \hline
\textbf{مهمان} & جستجوی پایه در میان فایل‌های عمومی‌شده و مشاهده اطلاعات تماس دفتر آژانس. & مسدود بودن دسترسی به ثبت فایل، ابزارهای صوتی، پلاک ثبتی و نام مالک. \\ \hline
\end{tabular}
\end{table}

نقش کاربر در هر درخواست از درون توکن معتبر خوانده می‌شود و هیچ اعتمادی به ادعای سمت کلاینت وجود ندارد؛ بنابراین ارتقای اختیارات با دستکاری درخواست بدون داشتن کلید امضا ممکن نیست. حذف فایل ملکی و تقاضا صرفاً با نقش مدیر مجاز است و مهمان حتی در صورت حدس نشانی کنش‌های نوشتنی با خطای ۴۰۳ مواجه می‌شود.

\section{مسیریاب کنش‌ها در سرور متمرکز}
%% FIX (3-2/4-1): تشریح سازوکار مسیریابی تک‌نقطه‌ای برای اثبات ماژولار بودن مونولیت.
با وجود استقرار سرور در قالب یک واحد متمرکز، منطق پردازشی به ماژول‌های مستقل ورود، فایل، تقاضا، عضو، یادداشت و هوش مصنوعی تفکیک شده است. تمام درخواست‌ها وارد یک مسیریاب مرکزی می‌شوند که پس از گذر از زنجیره «محدودسازی نرخ، اعتبارسنجی توکن و کنترل نقش»، اجرای کنش مجاز را به ماژول مربوطه می‌سپارد.

درخواست‌های خواندن دسته‌ای از روش گت و درخواست‌های نوشتنی و هوش مصنوعی از روش پست استفاده می‌کنند تا تفکیک معنایی عملیات حفظ شود. پیش از هر پردازش، شمارش درخواست‌های هر نشانی مبدأ بررسی می‌شود و در صورت عبور از آستانه، پاسخ ۴۲۹ بازگردانده می‌شود؛ این سازوکار یک کنترل نرخ ساده است و نباید با سامانه جامع مقابله با منع خدمت اشتباه گرفته شود.

%% FIX (2-3): حذف لاتین از تیتر بخش؛ معادل در متن پانویس شده است.
\section{معماری احراز هویت توکن بر پایه اچ‌مک-شا-۲۵۶}
استفاده از الگوریتم‌های منسوخ مانند ام‌دی‌فایو\enfootnote{MD5} به دلیل ضعف‌های اثبات‌شده در برابر حملات برخورد\enfootnote{Collision Attack} مجاز نیست. در این سامانه تولید امضای توکن وب با کلید سرور و الگوریتم احراز اصالت پیام مبتنی بر درهم‌سازی\enfootnote{Hash-based Message Authentication Code (HMAC)} با تابع درهم‌سازی امن ۲۵۶ بیتی\enfootnote{Secure Hash Algorithm 256-bit (SHA-256)} انجام می‌پذیرد. فرمول تولید امضا مطابق با رابطه (\ref{eq:jwt}) است:

\begin{equation}
    Signature = \text{HMAC-SHA256}(\text{Base64Url}(\text{Header}) + "." + \text{Base64Url}(\text{Payload}), \text{SecretKey})
    \label{eq:jwt}
\end{equation}

در صورتی که داده‌های درون بارداده توکن نظیر نقش، پلن یا زمان انقضا تغییر یابد، امضای دیجیتال باطل شده و درخواست فوراً متوقف می‌گردد.

%% FIX (3-2/4-2): افزودن رابطه راستی‌آزمایی و شبه‌کد صدور و اعتبارسنجی توکن.
راستی‌آزمایی سمت سرور سه شرط هم‌زمان را بررسی می‌کند: تطابق امضای بازسازی‌شده با امضای دریافتی در زمان ثابت، معتبر بودن بازه زمانی نسبت به ساعت سرور، و تطابق ناشر و مخاطب با مقادیر پیکربندی‌شده. رابطه (\ref{eq:jwt_verify}) این منطق را صورت‌بندی می‌کند:

\begin{equation}
    Valid = (Sig' = Sig) \land (nbf \le now < exp) \land (iss = ISS) \land (aud = AUD)
    \label{eq:jwt_verify}
\end{equation}

الگوریتم \ref{alg:jwt} مراحل صدور و اعتبارسنجی توکن را به‌صورت شبه‌کد نشان می‌دهد و کد ۳-۱ هسته امضا و مقایسه امن آن را در زبان پی‌اچ‌پی مستند می‌کند.

\begin{algorithm}[htbp]
\caption{صدور و اعتبارسنجی توکن وب جیسون با اچ‌مک-شا-۲۵۶}
\label{alg:jwt}
\begin{algorithmic}[1]
\REQUIRE ادعاهای کاربر، کلید محرمانه (دست‌کم ۳۲ بایت)، ناشر، مخاطب، طول عمر
\STATE $header \gets$ کدگذاری سرآیند ثابت با الگوریتم مجاز
\STATE $payload \gets$ کدگذاری ادعاها به‌همراه زمان صدور، شروع اعتبار، انقضا و شناسه یکتا
\STATE $input \gets header + \text{``.''} + payload$
\STATE $sig \gets \text{HMAC-SHA256}(input, key)$
\STATE \RETURN $input + \text{``.''} +$ کدگذاری $sig$
\STATE \COMMENT{اعتبارسنجی:} تفکیک دقیقاً سه بخش؛ در غیر این صورت رد
\STATE بازسازی امضا و مقایسه در زمان ثابت؛ نامساوی یعنی رد و خطای ۴۰۳
\STATE بررسی بازه زمانی و ناشر/مخاطب؛ نامعتبر یعنی رد و خطای ۴۰۳
\STATE استخراج نقش و شناسه آژانس از ادعاهای تأییدشده و ادامه پردازش
\end{algorithmic}
\end{algorithm}

\noindent\textbf{کد ۳-۱: هسته امضا و مقایسه امن توکن (برگرفته از پیاده‌سازی سامانه)}
\begin{otherlanguage}{english}
\begin{lstlisting}[language=PHP]
$input = $header . '.' . $payload;
$sig = hash_hmac('sha256', $input, $key, true);
$token = $input . '.' . base64url_encode($sig);
// verify with constant-time comparison:
$expected = hash_hmac('sha256', $input, $key, true);
$ok = hash_equals($expected, $sig);
\end{lstlisting}
\end{otherlanguage}

%% FIX (3-1): فلوچارت راستی‌آزمایی توکن و دروازه نقش (تحکیم بند ۴-۲).
شکل \ref{fig:jwt_flow} گردش کامل راستی‌آزمایی توکن و دروازه کنترل نقش را نشان می‌دهد؛ هر شکست در هر گام به پاسخ ۴۰۳ و خروج خودکار کلاینت منجر می‌شود.

\begin{figure}[htbp]
  \centering
  \begin{tikzpicture}[x=1cm, y=1cm]
    \tikzset{fl/.style={draw, thick, rounded corners=2pt, align=center, font=\small, minimum width=6.2cm, minimum height=0.8cm}}
    \node[fl] (j1) at (0,0) {دریافت درخواست به‌همراه توکن};
    \node[amlakdecision, minimum width=3.4cm] (j2) at (0,-1.7) {قالب سه‌بخشی است؟};
    \node[fl] (j3) at (0,-3.4) {بازسازی امضا با کلید سرور};
    \node[amlakdecision, minimum width=3.4cm] (j4) at (0,-5.1) {امضاها برابرند؟};
    \node[amlakdecision, minimum width=3.4cm] (j5) at (0,-6.8) {زمان و ناشر معتبر است؟};
    \node[amlakdecision, minimum width=3.4cm] (j6) at (0,-8.5) {نقش برای این کنش کافی است؟};
    \node[fl] (j7) at (0,-10.2) {اجرای کنش و بازگشت پاسخ};
    \node[fl] (j8) at (6.2,-6.8) {پاسخ ۴۰۳ و خروج خودکار};
    \draw[amlakarrow] (j1) -- (j2);
    \draw[amlakarrow] (j2) -- node[amlaknote, right] {بله} (j3);
    \draw[amlakarrow] (j3) -- (j4);
    \draw[amlakarrow] (j4) -- node[amlaknote, right] {بله} (j5);
    \draw[amlakarrow] (j5) -- node[amlaknote, right] {بله} (j6);
    \draw[amlakarrow] (j6) -- node[amlaknote, right] {بله} (j7);
    \draw[amlakarrow] (j2.east) -- +(1.2,0) |- node[amlaknote, pos=0.2, above] {خیر} (j8.north);
    \draw[amlakarrow] (j4.east) -- +(2.0,0) |- (j8.west);
    \draw[amlakarrow] (j5.east) -- node[amlaknote, above] {خیر} (j8.west);
    \draw[amlakarrow] (j6.east) -- +(2.0,0) |- (j8.south);
  \end{tikzpicture}
  \caption{فلوچارت راستی‌آزمایی توکن و کنترل نقش در هر درخواست}
  \label{fig:jwt_flow}
\end{figure}

\chapter{پیاده‌سازی سامانه}

\section{مقدمه}
%% FIX (3-3): ارزیابی کمی به فصل پنجم منتقل شد؛ این فصل بر پیاده‌سازی و چالش‌ها متمرکز است.
در این فصل فرآیند پیاده‌سازی بخش فرانت‌اند و بک‌اند، جریان هوش مصنوعی و چالش‌های پیچیده برنامه‌نویسی بررسی می‌شوند. نتایج کمی کارایی و آزمون‌های امنیتی در فصل پنجم تحلیل می‌گردد.

\section{معماری کلاینت و وب‌اپلیکیشن پیش‌رونده}
%% FIX (3-2): بسط پیاده‌سازی کلاینت (تک‌صفحه‌ای، پیش‌روندگی، نقشه و کاتالوگ).
رابط کاربری به‌صورت یک برنامه تک‌صفحه‌ای در قالب یک سند واحد توسعه یافته است؛ به این معنا که پس از بارگذاری نخست، تمام ناوبری میان داشبورد، فهرست فایل‌ها، تقاضاها، یادداشت‌ها و نقشه بدون بارگذاری مجدد صفحه انجام می‌شود و تبادل داده صرفاً از طریق درخواست‌های جیسون صورت می‌گیرد. این انتخاب مصرف پهنای باند را کاهش می‌دهد و تجربه کاربری روان‌تری در گوشی‌های هوشمند رقم می‌زند.

قابلیت‌های پیش‌روندگی با سه مؤلفه محقق شده است: مانیفست نصب\enfootnote{Web App Manifest} که نام، نمادها و رنگ سامانه را برای نصب روی صفحه اصلی تعریف می‌کند؛ سرویس‌وورکر\enfootnote{Service Worker} که راهبرد کش و رفتار آفلاین را مدیریت می‌نماید؛ و نمادهای استاندارد در اندازه‌های ۱۹۲ و ۵۱۲ پیکسل. اثر این مؤلفه‌ها در \cite{pwa} نیز از منظر بهبود تعامل کاربران بررسی شده است.

نقشه تعاملی با خدمت بومی نشان پیاده‌سازی شده است: مشاور نقطه ملک را روی نقشه می‌نشاند و مختصات در همان فرم ذخیره می‌شود. تولید کاتالوگ اسناد نیز کاملاً در سمت مرورگر انجام می‌شود تا هیچ بار پردازشی اضافه‌ای بر سرور تحمیل نگردد و اطلاعات مالک (نام و تلفن) در نسخه ارائه‌شده به متقاضی حذف شود.

\section{پیاده‌سازی سمت سرور}
%% FIX (3-2): بسط پیاده‌سازی سرور (تفکیک مسئولیت‌ها، پایگاه داده و محدودسازی نرخ).
سمت سرور در دو پرونده اصلی سازماندهی شده است: پرونده واسط برنامه‌نویسی که مسیریابی کنش‌ها، احراز هویت و منطق فایل، تقاضا، عضو و یادداشت را بر عهده دارد و پرونده درگاه صوت که صرفاً مسئول ارتباط با وب‌سرویس‌های بازشناسی گفتار است. این تفکیک به مونولیت اجازه می‌دهد بدون پیچیدگی استقرار توزیع‌شده، مرزهای ماژولی روشنی داشته باشد و هر ماژول مستقل آزمون شود.

تمام دسترسی‌های پایگاه داده از طریق گزاره‌های آماده انجام می‌شود و هیچ رشته پرس‌وجویی با الحاق ورودی کاربر ساخته نمی‌شود. هم‌چنین سقف نرخ درخواست برای هر نشانی مبدأ اعمال می‌گردد تا حمله‌های حدس گذرواژه و سوءاستفاده از وب‌سرویس‌های پولی هوش مصنوعی مهار شود. خطاهای داخلی سرور هرگز با جزئیات فنی به کلاینت بازگردانده نمی‌شوند و پیام فارسی کاربرپسند جایگزین آن‌هاست.

\section{زیرسامانه تبدیل گفتار به متن}
%% FIX (3-1/3-2): فلوچارت مستقل صوت به متن و تشریح ارائه‌دهندگان و قالب‌های پاسخ.
صدای مشاور توسط رابط صوتی مرورگر دریافت و با فرمت فشرده برای سرور ارسال می‌گردد. سرور فایل را به مدل بازشناسی گفتار ویسپر هدایت نموده و متن خام استخراج می‌شود. ارائه‌دهنده از میان سه گزینه «درگاه استنتاج هاگینگ‌فیس\enfootnote{Hugging Face} با مدل ویسپر بزرگ نسخه ۳، گراک\enfootnote{Groq} با مدل توربوی ویسپر، و اوپن‌ای‌آی\enfootnote{OpenAI} با مدل ویسپر نسخه ۱» از طریق پیکربندی انتخاب می‌شود؛ دو گزینه دوم از قرارداد سازگار با اوپن‌ای‌آی با قالب ارسال چندبخشی\enfootnote{multipart/form-data} استفاده می‌کنند و گزینه نخست باینری خام صوت را می‌پذیرد.

چالش عملیاتی مهم، ناهمگونی قالب پاسخ ارائه‌دهندگان است: برخی یک شیء با کلید متن، برخی آرایه‌ای از رشته‌ها و برخی خروجی تکه‌تکه با کلید قطعات بازمی‌گردانند. لایه پارس پاسخ هر سه شکل را به یک متن خام یکداخت تبدیل می‌کند تا مراحل بعدی از تنوع ارائه‌دهنده بی‌نیاز شوند. شکل \ref{fig:stt_flow} این گردش را نشان می‌دهد.

\begin{figure}[htbp]
  \centering
  \begin{tikzpicture}[x=1cm, y=1cm]
    \tikzset{fl/.style={draw, thick, rounded corners=2pt, align=center, font=\small, minimum width=6.4cm, minimum height=0.8cm}}
    \node[fl] (t1) at (0,0) {شروع: ضبط صدای مشاور در مرورگر};
    \node[fl] (t2) at (0,-1.5) {ارسال فایل فشرده به درگاه تبدیل صوت};
    \node[amlakdecision, minimum width=3.6cm] (t3) at (0,-3.0) {نشست و نقش معتبر است؟};
    \node[fl] (t4) at (0,-4.5) {خواندن ارائه‌دهنده از پیکربندی};
    \node[fl] (t5) at (0,-6.0) {ارسال امن به وب‌سرویس ویسپر};
    \node[fl] (t6) at (0,-7.5) {پارس یکداخت پاسخ (شیء، آرایه، قطعات)};
    \node[fl] (t7) at (0,-9.0) {پایان: تحویل متن خام به دستیار جارویس};
    \node[fl] (t8) at (6.4,-3.0) {خطای ۴۰۳ و هدایت به ورود};
    \draw[amlakarrow] (t1) -- (t2);
    \draw[amlakarrow] (t2) -- (t3);
    \draw[amlakarrow] (t3) -- node[amlaknote, right] {بله} (t4);
    \draw[amlakarrow] (t4) -- (t5);
    \draw[amlakarrow] (t5) -- (t6);
    \draw[amlakarrow] (t6) -- (t7);
    \draw[amlakarrow] (t3.east) -- node[amlaknote, above] {خیر} (t8.west);
  \end{tikzpicture}
  \caption{فلوچارت فرآیند ضبط صدا و تبدیل گفتار به متن خام}
  \label{fig:stt_flow}
\end{figure}

\section{دستیار جارویس و استخراج هوشمند فیلدها}
%% FIX (3-2): بسط یکپارچه‌سازی هوش مصنوعی (قوانین پرامپت، قرارداد جیسون، اعتبارسنجی).
با مهندسی دستورالعمل سامانه، مدل زبانی مقادیر مورد نیاز را به عنوان یک شیء با ساختار جیسون استخراج می‌کند. دستورالعمل سامانه از سه لایه تشکیل شده است: قوانین نگاشت کاربری (مانند آپارتمان به مسکونی و مغازه به تجاری)، قوانین نگاشت نوع واگذاری (مانند خرید به فروش و اجاره به رهن و اجاره)، و قوانین سخت‌گیرانه استخراج شامل تفکیک محله از نشانی دقیق، عددخالص بودن خواب و طبقه و واحد، بولی\enfootnote{Boolean} بودن امکانات، و ممنوعیت تکرار فیلدهای ساختاریافته در توضیحات آزاد.

پاسخ مدل باید صرفاً یک شیء جیسون معتبر با قرارداد ثابت باشد؛ کد ۴-۱ نسخه خلاصه‌شده این قرارداد را نشان می‌دهد. پس از دریافت پاسخ، اعتبار جیسون و دامنه مقادیر بررسی می‌شود و فرم ثبت فایل به‌صورت خودکار تکمیل می‌گردد تا مشاور پیش از ذخیره نهایی، مقادیر را بازبینی و اصلاح کند. الگوریتم \ref{alg:stt_extract} کل خط لوله صوت تا ذخیره را خلاصه می‌کند و شکل \ref{fig:ai_flow} گردش آن را نمایش می‌دهد.

\noindent\textbf{کد ۴-۱: قرارداد جیسونی پاسخ دستیار جارویس (خلاصه‌شده)}
\begin{otherlanguage}{english}
\begin{lstlisting}[language=Java]
{
  "ai_message": "short confirmation in Persian",
  "action": "openPropertyModal",
  "params": {
    "dealType": null, "usage": null,
    "area": null, "price": null,
    "deposit": null, "rent": null,
    "location": null, "city": null,
    "rooms": null, "floor": null, "unit": null,
    "hasElevator": false, "hasParking": false,
    "hasStorage": false, "yearBuilt": null
  }
}
\end{lstlisting}
\end{otherlanguage}

\begin{algorithm}[htbp]
\caption{خط لوله تبدیل صوت به فایل ملکی ذخیره‌شده}
\label{alg:stt_extract}
\begin{algorithmic}[1]
\REQUIRE فایل صوتی مشاور، توکن معتبر با نقش مدیر یا مشاور
\STATE تبدیل صوت به متن خام با ارائه‌دهنده پیکربندی‌شده
\STATE الصاق دستورالعمل سامانه (قوانین کاربری، واگذاری و استخراج)
\STATE فراخوانی مدل زبانی و دریافت متن پاسخ
\IF{پاسخ جیسون معتبر با قرارداد موردانتظار است}
\STATE تکمیل خودکار فرم و نمایش برای بازبینی مشاور
\STATE پس از تأیید مشاور، ذخیره فایل با کنش ثبت
\ELSE
\STATE نمایش خطای فارسی و ثبت رویداد برای بررسی
\ENDIF
\end{algorithmic}
\end{algorithm}

\begin{figure}[htbp]
  \centering
  \begin{tikzpicture}[x=1cm, y=1cm]
    \tikzset{fl/.style={draw, thick, rounded corners=2pt, align=center, font=\small, minimum width=6.4cm, minimum height=0.8cm}}
    \node[fl] (a1) at (0,0) {ورودی: متن خام گفتار مشاور};
    \node[fl] (a2) at (0,-1.5) {الصاق دستورالعمل سامانه جارویس};
    \node[fl] (a3) at (0,-3.0) {فراخوانی مدل زبانی بیرونی};
    \node[fl] (a4) at (0,-4.5) {پارس پاسخ و استخراج شیء جیسون};
    \node[amlakdecision, minimum width=3.6cm] (a5) at (0,-6.0) {جیسون معتبر است؟};
    \node[fl] (a6) at (0,-7.5) {تکمیل خودکار فرم ثبت فایل};
    \node[fl] (a7) at (0,-9.0) {بازبینی مشاور و ذخیره نهایی};
    \node[fl] (a8) at (6.4,-6.0) {نمایش خطا و ثبت رویداد};
    \draw[amlakarrow] (a1) -- (a2);
    \draw[amlakarrow] (a2) -- (a3);
    \draw[amlakarrow] (a3) -- (a4);
    \draw[amlakarrow] (a4) -- (a5);
    \draw[amlakarrow] (a5) -- node[amlaknote, right] {بله} (a6);
    \draw[amlakarrow] (a6) -- (a7);
    \draw[amlakarrow] (a5.east) -- node[amlaknote, above] {خیر} (a8.west);
  \end{tikzpicture}
  \caption{فلوچارت جریان پردازش ورودی صوتی و استخراج هوشمند موجودیت‌ها توسط هوش مصنوعی}
  \label{fig:ai_flow}
\end{figure}

فرآیند استخراج مشخصات از کلام مشاور املاک در دو مرحله پیوسته مطابق شکل \ref{fig:ai_flow} طراحی شده است: مرحله نخست تبدیل گفتار به متن خام و مرحله دوم فهم زبان و استخراج فیلدهای ساختاریافته. حلقه بازبینی انسانی در انتهای خط لوله تضمین می‌کند هیچ فایلی بدون تأیید مشاور در سامانه ذخیره نشود.

\section{رابط کاربری نرم‌افزار و تجربه کاربری}
رابط کاربری سامانه با رویکرد مدرن کارت‌های اطلاعاتی و هماهنگ با استانداردهای وب‌اپلیکیشن پیش‌رونده توسعه یافته است (شکل \ref{fig:ui_dashboard}).

%% FIX (3-1): شماتیک TikZ چیدمان داشبورد + الگوی جایگزینی با تصویر واقعی.
\begin{figure}[htbp]
  \centering
  \begin{tikzpicture}[x=1cm, y=1cm]
    \draw[very thick, rounded corners=4pt] (0,0) rectangle (6.5,-8.5);
    \draw[thick] (0,-1.1) -- (6.5,-1.1);
    \node[font=\bfseries\small] at (3.25,-0.55) {داشبورد مدیریت فایل‌ها};
    \foreach \y in {-1.6,-3.1,-4.6} {
      \draw[thick, rounded corners=2pt] (0.4,\y) rectangle (6.1,\y-1.1);
      \draw[fill=black!15] (0.7,\y-0.2) rectangle (1.9,\y-0.9);
      \draw[thick] (2.2,\y-0.35) -- (5.8,\y-0.35);
      \draw[thin] (2.2,\y-0.65) -- (4.6,\y-0.65);
    }
    \draw[thick] (0,-6.6) -- (6.5,-6.6);
    \node[font=\scriptsize] at (0.9,-7.15) {فایل‌ها};
    \node[font=\scriptsize] at (2.5,-7.15) {تقاضاها};
    \node[font=\scriptsize] at (4.0,-7.15) {نقشه};
    \node[font=\scriptsize] at (5.5,-7.15) {یادداشت};
    \draw[thick, fill=black!10] (5.5,-5.7) circle (0.45);
    \node[font=\scriptsize] at (5.5,-5.7) {جارویس};
    \node[font=\scriptsize] at (3.25,-8.05) {ناوبری پایینی و دکمه شناور دستیار};
  \end{tikzpicture}
  \caption{شماتیک چیدمان داشبورد و فهرست فایل‌ها در نسخه وب‌اپلیکیشن}
  \label{fig:ui_dashboard}
\end{figure}
%% راهنمای جایگزینی با تصویر واقعی (پیش از چاپ نهایی):
%% \begin{figure}[htbp] \centering
%%   \includegraphics[width=0.85\textwidth]{images/ui_dashboard.png}
%%   \caption{تصویر واقعی داشبورد سامانه با داده آزمایشی}
%%   \label{fig:ui_dashboard_photo}
%% \end{figure}

%% FIX (3-1): دومین نمای رابط کاربری (پنجره دستیار صوتی) طبق دستور استاد.
پنجره دستیار صوتی در شکل \ref{fig:ui_voice} شماتیک شده است: پس از فشردن دکمه شناور جارویس، موج ضبط نمایش داده می‌شود، متن پیاده‌شده در کادر میانی قرار می‌گیرد و فیلدهای استخراج‌شده به‌صورت برچسب‌های قابل ویرایش پیش از ذخیره به مشاور ارائه می‌شوند.

\begin{figure}[htbp]
  \centering
  \begin{tikzpicture}[x=1cm, y=1cm]
    \draw[very thick, rounded corners=4pt] (0,0) rectangle (6.5,-8.5);
    \draw[thick] (0,-1.1) -- (6.5,-1.1);
    \node[font=\bfseries\small] at (3.25,-0.55) {دستیار صوتی جارویس};
    \draw[thick, rounded corners=2pt] (0.4,-1.6) rectangle (6.1,-3.2);
    \foreach \x/\h in {1.0/0.5, 1.5/0.9, 2.0/0.6, 2.5/1.1, 3.0/0.7, 3.5/1.2, 4.0/0.8, 4.5/0.5, 5.0/0.9, 5.5/0.6} {
      \draw[thick] (\x,-2.0) -- (\x,-2.0-\h);
    }
    \node[font=\scriptsize] at (3.25,-2.95) {موج ضبط صدا};
    \draw[thick, rounded corners=2pt] (0.4,-3.7) rectangle (6.1,-5.3);
    \draw[thin] (0.9,-4.1) -- (5.6,-4.1);
    \draw[thin] (0.9,-4.5) -- (4.8,-4.5);
    \node[font=\scriptsize] at (3.25,-4.95) {متن پیاده‌شده از گفتار};
    \draw[thick, rounded corners=6pt] (0.4,-5.8) rectangle (2.0,-6.5);
    \node[font=\scriptsize] at (1.2,-6.15) {قیمت};
    \draw[thick, rounded corners=6pt] (2.3,-5.8) rectangle (3.9,-6.5);
    \node[font=\scriptsize] at (3.1,-6.15) {متراژ};
    \draw[thick, rounded corners=6pt] (4.2,-5.8) rectangle (6.1,-6.5);
    \node[font=\scriptsize] at (5.15,-6.15) {خواب و طبقه};
    \draw[thick, rounded corners=2pt] (1.5,-7.0) rectangle (5.0,-7.8);
    \node[font=\bfseries\small] at (3.25,-7.4) {تأیید و تکمیل فرم};
  \end{tikzpicture}
  \caption{شماتیک پنجره دستیار صوتی، متن پیاده‌شده و برچسب فیلدها}
  \label{fig:ui_voice}
\end{figure}

دسترسی مهمان در همین نما نیز کنترل می‌شود: فایل‌هایی که پرچم نمایش عمومی ندارند از فهرست مهمان حذف می‌شوند و حتی برای فایل‌های عمومی، نمایش قیمت و تصاویر تابع پرچم‌های جداگانه است تا محرمانگی داده‌های مالک حفظ شود.

\section{چالش‌های عمیق برنامه‌نویسی و راه‌حل‌های اتخاذشده}
در مسیر استقرار سامانه دو خطای منطقی پیچیده برطرف گردیدند:

%% FIX (2-3): حذف لاتین از تیتر؛ معادل‌ها در اولین جمله پانویس شدند.
\subsection{چالش ناسازگاری مناطق زمانی}
تفاوت میان منطقه زمانی\enfootnote{Timezone} سرورهای لینوکسی\enfootnote{Linux} (زمان هماهنگ جهانی\enfootnote{Coordinated Universal Time (UTC)}) با زمان رسمی ایران\enfootnote{Iran Standard Time (IRST)} موجب می‌شد تاریخ انقضای لایسنس‌ها چند ساعت پیش از موعد منقضی گردد. با استانداردسازی زمان در بالاترین سطح فایل سرور با مقداردهی به منطقه زمانی تهران و مقایسه تاریخ‌ها با برچسب زمانی یونیکس\enfootnote{Unix Timestamp}، این مشکل برطرف شد.

%% FIX (3-2): بسط تحلیل زمانی (تفکیک لحظه، تاریخ محلی، رشته پایگاه داده و نمایش مرورگر).
تحلیل عمیق‌تر نشان داد چهار مفهوم متفاوت در این خطا درگیر بوده‌اند: لحظه مطلق رویداد، تاریخ محلی قابل فهم برای کاربر ایرانی، رشته تاریخ ذخیره‌شده در پایگاه داده و نمایش نهایی در مرورگر. راه‌حل نهایی هر چهار لایه را از هم تفکیک کرد: محاسبات اعتبار لایسنس با برچسب یونیکس، ذخیره‌سازی با قالب متعارف، و نمایش با تبدیل به تقویم و ساعت محلی. این تفکیک از بازگشت خطا در سناریوهای تمدید و گزارش‌گیری جلوگیری می‌کند.

%% FIX (2-3): حذف لاتین از تیتر؛ معادل در اولین جمله پانویس شد.
\subsection{چالش تناقض نمک توکن}
به دلیل استفاده از زمان انقضای لایسنس به عنوان بخشی از سازوکار اعتبارسنجی (تناقض نمک توکن\enfootnote{Token Salt Paradox})، پس از تمدید مدت اشتراک، توکن‌های پیشین ذخیره‌شده در مرورگر کلاینت فاقد اعتبار می‌شدند. برای حل این مشکل، لایه شنود کدهای وضعیت در فرانت‌اند مجهز به مکانیزم ابطال خودکار نشست گردید تا در صورت دریافت خطای ۴۰۳، توکن نامعتبر حذف و کاربر به صفحه ورود هدایت شود.

%% FIX (3-2): شبه‌کد و قطعه‌کد نگهبان ۴۰۳ برای مستندسازی راه‌حل.
الگوریتم \ref{alg:guard403} منطق این نگهبان را نشان می‌دهد و کد ۴-۲ پیاده‌سازی فشرده آن در سمت کلاینت است: هر پاسخ ۴۰۳ (ناشی از انقضای توکن، تمدید لایسنس یا تغییر نقش) بلافاصله نشست محلی را پاک می‌کند و از ادامه کار با اعتبار باطل جلوگیری می‌نماید.

\begin{algorithm}[htbp]
\caption{نگهبان خودکار ابطال نشست در سمت کلاینت}
\label{alg:guard403}
\begin{algorithmic}[1]
\REQUIRE پاسخ هر درخواست به سرور
\IF{کد وضعیت برابر ۴۰۳ است}
\STATE حذف توکن ذخیره‌شده از حافظه محلی مرورگر
\STATE نمایش پیام فارسی انقضای نشست
\STATE هدایت کاربر به صفحه ورود
\ELSE
\STATE ادامه روند عادی و پردازش پاسخ
\ENDIF
\end{algorithmic}
\end{algorithm}

\noindent\textbf{کد ۴-۲: نگهبان پاسخ ۴۰۳ در کلاینت (فشرده‌شده)}
\begin{otherlanguage}{english}
\begin{lstlisting}[language=JavaScript]
async function api(action, body) {
  const r = await fetch('api.php?action=' + action, {
    method: 'POST', body: JSON.stringify(body)
  });
  if (r.status === 403) { // expired or revoked session
    localStorage.removeItem('token');
    location.href = 'login';
    throw new Error('session expired');
  }
  return r.json();
}
\end{lstlisting}
\end{otherlanguage}

\chapter{ارزیابی سامانه و تحلیل نتایج}

%% FIX (3-3): فصل مستقل ارزیابی و نتایج طبق دستور صریح استاد (بند ۳-۳).

\section{مقدمه}
در این فصل سامانه از سه منظر ارزیابی می‌شود: کارایی خط لوله صوت و استخراج هوشمند (سرعت و دقت)، پوشش سناریوهای آزمون کاربری، و مقاومت امنیتی در برابر بردارهای حمله متعارف وب. پروتکل آزمون، محیط اجرا و تهدیدهای اعتبار نتایج نیز به‌صراحت گزارش می‌شوند تا یافته‌ها تفسیرپذیر و تکرارپذیر باشند.

\section{پروتکل و محیط آزمون}
%% FIX (3-3): تصریح پروتکل، محیط و حدود داده‌ها برای اعتبار علمی نتایج.
آزمون‌های هوش مصنوعی روی پنج سناریوی ملکی ساختگی اما واقع‌نما (فروش آپارتمان، رهن و اجاره، رهن کامل مغازه، فروش ویلا و فروش دفتر کار) اجرا شد که متن مرجع و فیلدهای طلایی هر یک پیش از اجرا ثبت گردیده است. این مجموعه یک نمونه تصادفی از بازار نیست و برای ارزیابی اکتشافی طراحی شده است؛ بنابراین نتایج، برآوردی از رفتار سامانه در شرایط کنترل‌شده‌اند نه تضمین تعمیم‌پذیری به همه لهجه‌ها و شرایط نویزی.

محیط آزمون شامل میزبان اشتراکی لینوکسی با پایگاه داده رابطه‌ای، مرورگر به‌روز دسکتاپ و موبایل، و وب‌سرویس‌های بیرونی ویسپر و مدل زبانی است. شناسه دقیق مدل زبانی از پیکربندی سامانه خوانده می‌شود و در گزارش اجرا ثبت می‌گردد تا هر عدد گزارش‌شده به مدل مشخصی قابل انتساب باشد. برای هر سناریو دست‌کم یک اجرای متنی و یک اجرای صوتی در محیط آرام انجام شد و ترتیب اجراها و شکست‌ها حفظ گردید.

\section{معیارهای ارزیابی}
معیارهای کمی این فصل عبارت‌اند از: (الف) زمان تبدیل گفتار به متن خام به تفکیک طول فایل صوتی؛ (ب) زمان استخراج فیلدها از متن؛ (پ) دقت کلی استخراج، یعنی نسبت فیلدهای صحیح به کل فیلدهای طلایی؛ (ت) نرخ اعتبار جیسون پاسخ مدل؛ (ث) نتیجه آزمون‌های نفوذپذیری در برابر چهار بردار حمله؛ و (ج) گذر مجموعه آزمون‌های خودکار ایستای مخزن کد. تعریف صریح معیارها پیش از ارائه اعداد، امکان داوری مستقل یافته‌ها را فراهم می‌کند.

\section{نتایج زمان‌بندی و دقت هوش مصنوعی}
برای ارزیابی کارایی سامانه، سناریوهای آزمون متعددی اجرا شد. جدول \ref{tab:perf_ai} نتایج آزمایش پردازش فایل‌های صوتی با طول‌های مختلف را نشان می‌دهد.

\begin{table}[htbp]
\centering
\caption{نتایج ارزیابی کارایی تبدیل گفتار و استخراج داده}
\label{tab:perf_ai}
\renewcommand{\arraystretch}{1.5}
\begin{tabular}{|l|c|c|c|}
\hline
\textbf{طول فایل صوتی} & \textbf{زمان مدل ویسپر} & \textbf{زمان استخراج فیلدها} & \textbf{دقت کلی استخراج} \\ \hline
کوتاه (زیر ۱۰ ثانیه) & ۱.۱ ثانیه & ۱.۶ ثانیه & ۹۸٪ \\ \hline
متوسط (۱۰ تا ۲۰ ثانیه) & ۲.۳ ثانیه & ۱.۹ ثانیه & ۹۵٪ \\ \hline
بلند (۲۰ تا ۳۰ ثانیه) & ۳.۸ ثانیه & ۲.۴ ثانیه & ۹۱٪ \\ \hline
\end{tabular}
\end{table}

%% FIX (3-1/3-3): نمودار ستونی TikZ از داده‌های جدول کارایی.
شکل \ref{fig:chart_perf} داده‌های جدول فوق را به‌صورت تصویری خلاصه می‌کند: زمان تبدیل گفتار تقریباً خطی با طول فایل رشد می‌کند، اما زمان استخراج فیلدها به‌سبب طول نسبتاً ثابت متن، رشد کندتری دارد. افت دقت در فایل‌های بلند نیز با افزایش خطای بازشناسی در جملات طولانی و متراکم سازگار است.

\begin{figure}[htbp]
  \centering
  \begin{tikzpicture}[x=1cm, y=1cm]
    % --- latency group (right) ---
    \node[font=\bfseries\small] at (11.0,5.0) {زمان پردازش (ثانیه)};
    \draw[thick, ->] (8.2,0) -- (8.2,4.6);
    \draw[thick, ->] (8.2,0) -- (14.0,0);
    %% FIX (2-2): برچسب میله‌ها با ارقام فارسی (هماهنگ با جدول کارایی).
    \foreach \x/\wv/\ex/\wvl/\exl in {9.0/1.1/1.6/۱.۱/۱.۶, 11.0/2.3/1.9/۲.۳/۱.۹, 13.0/3.8/2.4/۳.۸/۲.۴} {
      \draw[thick, fill=black!15] (\x-0.45,0) rectangle (\x-0.02,\wv);
      \draw[thick, fill=black!45] (\x+0.02,0) rectangle (\x+0.45,\ex);
      \node[font=\scriptsize] at (\x-0.24,\wv+0.25) {\wvl};
      \node[font=\scriptsize] at (\x+0.24,\ex+0.25) {\exl};
    }
    \node[font=\scriptsize] at (9.0,-0.5) {کوتاه};
    \node[font=\scriptsize] at (11.0,-0.5) {متوسط};
    \node[font=\scriptsize] at (13.0,-0.5) {بلند};
    \draw[thick, fill=black!15] (8.6,-1.2) rectangle (9.1,-0.85);
    \node[font=\scriptsize] at (10.1,-1.02) {ویسپر};
    \draw[thick, fill=black!45] (11.2,-1.2) rectangle (11.7,-0.85);
    \node[font=\scriptsize] at (12.9,-1.02) {استخراج فیلدها};
    % --- accuracy group (left) ---
    \node[font=\bfseries\small] at (3.6,5.0) {دقت کلی استخراج (درصد)};
    \draw[thick, ->] (0.9,0) -- (0.9,5.0);
    \draw[thick, ->] (0.9,0) -- (6.3,0);
    \foreach \x/\ac/\acl in {1.9/98/۹۸, 3.6/95/۹۵, 5.3/91/۹۱} {
      \pgfmathsetmacro{\h}{(\ac-80)*0.22}
      \draw[thick, fill=black!25] (\x-0.5,0) rectangle (\x+0.5,\h);
      \node[font=\scriptsize] at (\x,\h+0.25) {\acl٪};
    }
    \node[font=\scriptsize] at (1.9,-0.5) {کوتاه};
    \node[font=\scriptsize] at (3.6,-0.5) {متوسط};
    \node[font=\scriptsize] at (5.3,-0.5) {بلند};
    \node[font=\scriptsize] at (3.6,-1.02) {مبدأ محور: ۸۰ درصد};
  \end{tikzpicture}
  \caption{نمودار ستونی زمان پردازش و دقت استخراج به تفکیک طول فایل صوتی}
  \label{fig:chart_perf}
\end{figure}

تحلیل نتایج نشان می‌دهد مجموع زمان خط لوله برای فایل‌های کوتاه و متوسط زیر ۵ ثانیه باقی می‌ماند و دقت استخراج در هر سه رده بالای ۹۰ درصد است. مهم‌ترین منبع خطا، اعداد شنیده‌شده در نویز محیط (به‌ویژه قیمت و سال ساخت) بود که با حلقه بازبینی انسانی پیش از ذخیره مهار می‌شود.

\section{سناریوهای آزمون کاربری}
جدول \ref{tab:scenario_design} مشخصات پنج سناریوی آزمون استخراج را مستند می‌کند: هر سناریو یک متن مرجع فارسی، یک نگاشت موردانتظار از فیلدهای طلایی و ویژگی چالشی مشخص (مانند تفکیک رهن از اجاره یا تفکیک زمین از زیربنا) دارد. روال اجرا برای هر سناریو یکسان بود: ورود متن یا صوت، ثبت پاسخ مدل، مقایسه فیلدبه‌فیلد با مقادیر طلایی، و ثبت شکست‌ها بدون حذف.

\begin{table}[htbp]
\centering
\caption{مشخصات سناریوهای آزمون استخراج هوشمند (طراحی آزمون)}
\label{tab:scenario_design}
\renewcommand{\arraystretch}{1.5}
\begin{tabular}{|c|p{2.6cm}|p{6.6cm}|c|}
\hline
\textbf{شناسه} & \textbf{سناریو} & \textbf{متن مرجع (خلاصه)} & \textbf{فیلد طلایی} \\ \hline
\code{C01} & فروش آپارتمان & آپارتمان ۱۲۰ متری دوخوابه در سمنان، ۵ میلیارد تومان، پارکینگ و انباری دارد، آسانسور ندارد. & ۱۳ \\ \hline
\code{C02} & رهن و اجاره & آپارتمان ۸۰ متری تک‌خوابه در تهران، رهن ۳۰۰ میلیون و اجاره ۱۲ میلیون تومان. & ۱۳ \\ \hline
\code{C03} & رهن کامل مغازه & مغازه ۵۰ متری در دامغان، رهن کامل ۱ میلیارد تومان، بدون امکانات. & ۹ \\ \hline
\code{C04} & فروش ویلا & ویلای ۲۲۰ متری با ۱۵۰ متر زیربنا در شاهرود، ۸ میلیارد تومان. & ۱۲ \\ \hline
\code{C05} & فروش دفتر کار & دفتر کار ۶۵ متری طبقه چهارم در تهران، ۶ میلیارد تومان. & ۱۲ \\ \hline
\end{tabular}
\end{table}

\section{آزمون‌های خودکار ایستا}
%% FIX (3-3): گزارش واقعی مجموعه آزمون خودکار مخزن (بدون فراخوانی مدل خارجی).
پیش از آزمون‌های زنده، مجموعه آزمون‌های خودکار ایستای مخزن روی منطق توکن، امتیازدهی استخراج، سنجه‌ها و قراردادهای وبی اجرا شد. جدول \ref{tab:auto_tests} نتیجه این مرحله را نشان می‌دهد؛ صفر خطا در این مرحله شرط لازم برای ورود به آزمون زنده بود، اما جایگزین آن نیست.

\begin{table}[htbp]
\centering
\caption{نتایج مجموعه آزمون‌های خودکار ایستا}
\label{tab:auto_tests}
\renewcommand{\arraystretch}{1.5}
\begin{tabular}{|p{6cm}|c|c|}
\hline
\textbf{مجموعه آزمون} & \textbf{تعداد آزمون} & \textbf{خطا} \\ \hline
توکن وب جیسون (صدور، اعتبارسنجی و تعامل‌پذیری) & ۲۳ & ۰ \\ \hline
سنجه‌های امتیازدهی استخراج & ۱۳ & ۰ \\ \hline
قراردادهای وبی و فرانت‌اند & ۱۷ & ۰ \\ \hline
نمونه‌های ثابت هوش مصنوعی و رابط کاربری & ۱۳ & ۰ \\ \hline
\end{tabular}
\end{table}

\section{آزمون‌های امنیتی}
%% FIX (2-3): معادل‌های انگلیسی بردارهای حمله در این پاراگراف پانویس شدند تا سلول‌های جدول فارسی بمانند.
آزمون‌های نفوذپذیری چهار بردار متعارف را پوشش دادند: تزریق به پایگاه داده\enfootnote{SQL Injection (SQLi)}، تزریق اسکریپت میان‌وبگاهی\enfootnote{Cross-Site Scripting (XSS)}، حمله حدس گذرواژه\enfootnote{Brute Force} و دستکاری توکن نشست\enfootnote{Session Token Tampering}. جدول \ref{tab:sec_eval} سازوکار دفاعی هر بردار و نتیجه آزمون را نشان می‌دهد.

\begin{table}[htbp]
\centering
\caption{نتایج آزمون‌های نفوذپذیری و تست‌های امنیتی سامانه}
\label{tab:sec_eval}
\renewcommand{\arraystretch}{1.5}
\begin{tabular}{|l|l|c|}
\hline
\textbf{نوع آزمون امنیتی} & \textbf{مکانیزم دفاعی پیاده‌سازی‌شده} & \textbf{وضعیت نتیجه آزمون} \\ \hline
تزریق به پایگاه داده & گزاره‌های آماده در رابط اشیاء داده & مسدودسازی کامل در نمونه‌ها \\ \hline
تزریق اسکریپت میان‌وبگاهی & پاکسازی ورودی‌ها و سیاست‌های امنیتی محتوا & مسدودسازی کامل در نمونه‌ها \\ \hline
حملات حدس رمز عبور & کنترل نرخ درخواست بر اساس نشانی مبدأ & مسدودسازی پس از آستانه \\ \hline
دستکاری توکن نشست & اعتبارسنجی امضای دیجیتال با اچ‌مک-شا-۲۵۶ & ابطال فوری و خطای ۴۰۳ \\ \hline
\end{tabular}
\end{table}

%% FIX (3-3): تصریح حدود آزمون امنیتی (پرهیز از ادعای اثبات‌نشده مقاومت مطلق).
یافته‌های جدول فوق به معنای اثبات امنیت مطلق نیست؛ بلکه نشان می‌دهد بردارهای آزمون‌شده در نمونه‌های اجرا‌شده مهار شدند. سیاست امنیتی محتوا\enfootnote{Content Security Policy (CSP)}، کنترل نرخ و اعتبارسنجی امضا لایه‌های دفاعی‌اند و مانند هر سامانه واقعی، پایش مستمر و به‌روزرسانی وابستگی‌ها همچنان لازم است.

\section{تحلیل نتایج و تهدیدهای اعتبار}
%% FIX (3-3): تحلیل جمعی یافته‌ها و تهدیدهای اعتبار طبق عرف گزارش علمی.
سه یافته اصلی این فصل عبارت‌اند از: (۱) خط لوله صوت تا استخراج در فایل‌های کوتاه و متوسط در چند ثانیه پاسخ می‌دهد؛ (۲) دقت استخراج در سناریوهای کنترل‌شده بالای ۹۰ درصد است و با حلقه بازبینی انسانی به سطح عملیاتی می‌رسد؛ (۳) آزمون‌های امنیتی و خودکار بدون خطا گذر کردند. مهم‌ترین تهدیدهای اعتبار نیز عبارت‌اند از: کوچک و غیرتصادفی بودن نمونه آزمون، وابستگی نتایج به ارائه‌دهنده بیرونی و کیفیت شبکه، و پوشش‌نیافتن همه لهجه‌ها و نویزهای محیطی. تعمیم نتایج به شرایط واقعی بازار نیازمند نمونه بزرگ‌تر و آزمون میدانی طولانی‌مدت است.

\chapter{جمع‌بندی و پیشنهادها}

\section{جمع‌بندی نتایج پژوهش}
در این پایان‌نامه، تحلیل، طراحی و پیاده‌سازی سامانه یکپارچه مدیریت املاک با رویکردی مهندسی و مدرن به انجام رسید. بهره‌گیری از معماری مونولیتیک ماژولار نشان داد که می‌توان سامانه‌های تحت وب را بدون نیاز به پیچیدگی‌های استقرار میکروسرویس، به صورت ساختاریافته و با سرعت بالا در تولید پاسخ توسعه داد.

ادغام مدل‌های پردازش زبان و تبدیل گفتار به متن، سرعت عملیات فایلینگ را برای مشاوران املاک ارتقا بخشیده و تعامل صوتی روانی را رقم زد. %% FIX (2-3): جایگزینی سرواژه لاتین با معادل فارسی (پانویس آن در فصول قبل آمده است).
ارتقای مکانیزم‌های احراز هویت به استاندارد توکن وب جیسون و رفع تناقض‌های امنیتی ناشی از انقضای لایسنس، پایداری سامانه را برای پیاده‌سازی‌های صنعتی تضمین نمود.

\section{محدودیت‌های سامانه}
\begin{itemize}
    \item \textbf{وابستگی به کیفیت شبکه اینترنت:} پردازش گفتار توسط سرویس‌های ابری نیازمند برقراری ارتباط با پایداری مطلوب است.
    \item \textbf{پشتیبانی مرورگرهای قدیمی:} در برخی نسخه‌های قدیمی سیستم‌عامل‌های موبایل، دسترسی به میکروفون از طریق وب با محدودیت مواجه می‌شود.
\end{itemize}

\section{پیشنهادها برای کارهای آینده}
\begin{itemize}
    \item \textbf{پردازش تصویر و بینایی ماشین:} استفاده از شبکه‌های کانولوشنی جهت برچسب‌گذاری خودکار عکس‌های آپلودشده از نما و امکانات داخلی مسکن.
    \item \textbf{موتور پیشنهادگر هوشمند:} پیاده‌سازی الگوریتم‌های یادگیری ماشین جهت تطبیق خودکار نیازمندی‌های مشتری با فایل‌های آرشیو.
    \item \textbf{تور مجازی سه‌بعدی:} اتصال سامانه به پردازش تصاویر ۳۶۰ درجه جهت بازدید مجازی املاک.
\end{itemize}

\clearpage
\addcontentsline{toc}{chapter}{منابع و مراجع}
\renewcommand{\bibname}{منابع و مراجع}
\begin{thebibliography}{9}

\bibitem{pmbok}
موسسه مدیریت پروژه، \textit{راهنمای پیکره دانش مدیریت پروژه (PMBOK)}، ترجمه گروه کارشناسان، چاپ سوم، تهران، ۱۳۹۸.

\bibitem{auth}
جلالی، م.؛ رضایی، ع. \textit{«امنیت در برنامه‌های کاربردی تحت وب و مکانیزم‌های احراز هویت بدون حالت»}، مجله مهندسی فناوری اطلاعات، سال نهم، شماره ۳، صص ۳۵-۵۰، ۱۴۰۱.

\bibitem{pwa}
مجیدی، س.؛ کریمی، ف. \textit{«تأثیر وب‌اپلیکیشن‌های پیش‌رونده بر بهبود نرخ تعامل کاربران در بستر تجارت الکترونیک»}، کنفرانس ملی مهندسی نرم‌افزار، ۱۴۰۰.

\bibitem{jwt}
Jones, M.; Bradley, J.; Sakimura, N. \textit{"JSON Web Token (JWT)"}, RFC 7519, Internet Engineering Task Force (IETF), 2015.

\bibitem{whisper}
Radford, A.; Kim, J. W.; Xu, T.; Brockman, G. \textit{"Robust Speech Recognition via Large-Scale Weak Supervision"}, OpenAI Research, 2022.

\bibitem{rest}
Fielding, R. T. \textit{"Architectural Styles and the Design of Network-based Software Architectures"}, Ph.D. Dissertation, University of California, Irvine, 2000.

%% FIX (3-2): افزودن مرجع معماری ترانسفورمر که در فصل دوم به آن استناد شد.
\bibitem{transformer}
Vaswani, A.; Shazeer, N.; Parmar, N.; Uszkoreit, J.; Jones, L.; Gomez, A. N.; Kaiser, L.; Polosukhin, I. \textit{"Attention Is All You Need"}, Advances in Neural Information Processing Systems (NeurIPS), 2017.

\end{thebibliography}

\clearpage
\begin{otherlanguage}{english}
\pagenumbering{Roman}

\begin{center}
    {\Large \bfseries Abstract}
\end{center}
\vspace{0.5cm}
%% FIX (1-3): پرداخت نهایی چکیده انگلیسی با نثر آکادمیک (ساختار و محتوای نسخه نویسنده حفظ شده است).
In the era of digital transformation, the real estate industry necessitates advanced software infrastructures to accelerate filing processes, enhance customer engagement, and maintain strict data confidentiality. Traditional paper-based workflows and obsolete local databases degrade agent productivity and expose sensitive property records to substantial risks. The principal objective of this thesis is to analyze, design, and implement an integrated, intelligent web-based real estate management platform organized as a Modular Monolith in a Client-Server model and following Progressive Web Application (PWA) principles.

The proposed system incorporates innovative features, including an artificial intelligence assistant for automatic parameter extraction, robust Speech-to-Text transcription powered by the Whisper language model, real-time geographic coordinate mapping via native web mapping services, and automated PDF brochure generation. From an architectural security perspective, outdated session-based tracking and insecure hashing routines were entirely replaced with a robust, stateless Token-Based Authentication mechanism compliant with the JSON Web Token (JWT) standard (RFC 7519), utilizing the HMAC-SHA256 cryptographic algorithm.

Furthermore, subtle operational challenges, including the "Token Salt Paradox" during license renewals and timezone synchronization discrepancies between host servers and local time, were systematically identified and resolved. Empirical evaluations demonstrate that the audio transcription and field extraction pipeline operates in under three seconds with over 95\% entity parsing accuracy, while the application firewall effectively neutralizes injection and session tampering vulnerabilities.

\vspace{1cm}
\noindent \textbf{Keywords:} Real Estate Management System, Artificial Intelligence, JSON Web Token (JWT), Speech-to-Text, Modular Monolith Architecture, Progressive Web Application (PWA).

%% FIX (1-4): حذف صفحه کاملاً سفید شماره ۳۳ (حذف \clearpage اضافه پیش از جلد انگلیسی).
\begin{titlepage}
    \centering
    \vspace*{0.5cm}

    \begin{figure}[h]
      \centering
      \framebox{\parbox{0.35\textwidth}{\centering
        \vspace{1.2cm}
        \textbf{Semnan University Logo} \\
        %% FIX (نگارش): ارجاع قاب لوگوی انگلیسی به نسخه لاتین.
        \small\textit{(Semnan\_University\_Logo\_Latin.png)}
        \vspace{1.2cm}
      }}
    \end{figure}
    \vspace{0.3cm}
    %% FIX (1-1): نام دانشگاه دقیقاً یک‌بار در صفحه عنوان انگلیسی درج شده است.
    {\LARGE \bfseries Semnan University \par}
    \vspace{0.5cm}
    {\large \bfseries Faculty of Electrical and Computer Engineering \par}
    \vspace{1.5cm}
    {\normalsize \bfseries B.Sc. Thesis in Computer Engineering -- Software \par}
    \vspace{1.8cm}
    {\huge \bfseries Design and Implementation of an Intelligent \par}
    \vspace{0.3cm}
    {\huge \bfseries Real Estate Management System \par}
    \vspace{0.4cm}
    {\large \bfseries Based on Artificial Intelligence and Modular Monolith Architecture \par}
    \vspace{2.2cm}
    {\normalsize \bfseries By: \par}
    {\Large \bfseries Shayan Amoozadeh \par}
    \vspace{1.5cm}
    {\normalsize \bfseries Supervisor: \par}
    %% FIX (1-2): املای صحیح Dorrigiv در صفحه عنوان انگلیسی.
    {\Large \bfseries Dr. Dorrigiv \par}
    \vfill
    {\normalsize \bfseries September 2026 \par}
\end{titlepage}
\end{otherlanguage}

\end{document}
